\section{Interoperability in \acl{bps}}
We assume that a set $\mathbb{P}$ of broadcast protocols is in execution in a deployed mobile network of $\mathbb{N}$ nodes where each node $n_{i} \in \mathbb{N}$ runs only one protocol $p_{j} \in \mathbb{P}$ at a time.
%
Our aim is letting a node $n_{i}$ chose in an autonomous way a certain protocol $p_{j}$, this decision is based in the local information available at the node.
%
Such information may be, and it is not limited to, remaining battery, signal strength of received packets or the node relative speed.
%
But we omit the case where $n_{i}$ is running more than one protocol concurrently.

%%
When $\mathbb{P}$ contains only one protocol the reachability of messages depends on the performance of such approach.
%
Notice that in our hybrid approach we do not intent to alter the internals of any protocol in $\mathbb{P}$.
%
On the contrarily, when the cardinality of $\mathbb{P}$ is $|\mathbb{P}| > 1$ packets lost may take place.

\begin{figure}[t!]
  \centering
  \includegraphics[scale=0.20]{packets-lost-ex}
  \caption{TBD}
  \label{packets-lost-ex}
\end{figure}

%
Figure~\ref{packets-lost-ex} depicts a situation where a set of nodes $\mathbb{T}=\{t_{i} \mid i \in \{1, \ldots, 6\}\}$ run a topology-based protocol while nodes in $\mathbb{F}=\{f_{j} \mid i \in \{1, \ldots, 3\}\}$ run blind flooding.
%
Broadcast sessions initiated at $t_{1}$ form a tree such that a forwarding set composed of $t_{1}, t_{2}, t_{3}$ guarantees the message reception of every node at $\mathbb{T}$.
%
Even if nodes in $t_{4}, t_{5}, t_{6}$ are within the communication range of some nodes in $\mathbb{F}$, their role as pure receivers in the tree avoids any node in $\mathbb{F}$ hearing from messages.

\subsection{Detection of border nodes}

Algorithm~\ref{interop-alg} is a simple mechanism we propose to deal with interoperability issues between nodes running different protocols.
%%%% INTEROPERABILITY ALGORITHM %%%%

\noindent
\SetInd{0.0em}{0.60em}
%\SetAlgoNoLine
\begin{algorithm}
\footnotesize
\DontPrintSemicolon
\caption{Interoperability mechanism to deal with more than one running protocol in a MANET.}
\label{interop-alg}
\vars{
	% VARIABLES
	node\_id \textit{: string} \tcp*[1]{node identifier}
	running\_alg \textit{: string} \tcp*[1]{running algorithm}
	timeout\textit {: double} \tcp*[1]{maximum waiting time to start the procedure of border dectection}
	am\_i\_border \textit{: boolean} \tcp*[1]{says wether the local node is a border node or not}

	% SETS
	known\_algs \textit{: set} \tcp*[1]{previously known foreign protocols}

	% MAPS
	timers \textit{: dictionary} \tcp*[1]{timers of border detection procedure}
	known\_nodes \textit{: dictionary} \tcp*[1]{previously known nodes from a foreign protocol}
}

\func{ \receive{msg} \label{interop-alg:rcv-fun} }{ 
	\eIf{ msg.sender\_alg = running\_alg \Ork am\_i\_border }{
		broadcast(msg) \; \label{interop-alg:bro}
	}{
		\Switch{msg}{
			\Case{ ControlMsg }{

				known\_nodes[msg.sender\_alg] $\gets$ msg.sender\_id \; \label{interop-alg:save-sender}
				
				\If{ msg.sender\_alg $\notin$ known\_algs }{
				
					known\_algs $\gets$ known\_algs $\cup$ msg.sender\_alg \;	\label{interop-alg:save-alg}
					delay $\gets$ $new$ Random(0, timeout) \;					
					timers[msg.sender\_alg] $\gets$ $new$ Timer(delay) \; \label{interop-alg:new-timer}
					
					\TimerExpires{

						b $\gets$ $new$ BorderMsg() \; \label{interop-alg:new-border-msg}
						b.chosen\_node $\gets$ known\_nodes[msg.sender\_alg].pickOne() \label{interop-alg:new-border-msg:chosen} \;
						b.fwd\_no $\gets$ MAX\_FWD\_NO \; \label{interop-alg:new-border-msg:hops}
						b.sender\_id $\gets$ node\_id \; \label{interop-alg:new-border-msg:sender}						
						b.sender\_alg $\gets$ running\_alg \; \label{interop-alg:new-border-msg:sender-alg}
						b.foster\_by $\gets$ msg.sender\_alg \; \label{interop-alg:new-border-msg:alg}
						forward(b) \; \label{interop-alg:fwd}

					}
				}
			}
			\Case{ BorderMsg }{
			
				msg.fwd\_no $\gets$ msg.fwd\_no - 1 \; \label{interop-alg:border-msg-case:dec}
				
				\If{ msg.fwd\_no $\neq$ 0 }{
					forward(msg) \; \label{interop-alg:border-msg-case:fwd}
				}
				
				\If{ running\_alg = msg.foster\_by \And msg.foster\_by $\notin$ known\_algs }{
				
					known\_algs $\gets$ known\_algs $\cup$ msg.foster\_by \; \label{interop-alg:border-msg-case:update-known-algs-1}

					\If{ node\_id = msg.chosen\_node }{					
						am\_i\_border $\gets$ True \; \label{interop-alg:border-msg-case:select-as-border}
					}
					
				}
				
				\If{ running\_alg = msg.sender\_alg \And msg.sender\_alg $\notin$ known\_algs }{
				
					known\_algs $\gets$ known\_algs $\cup$ msg.sender\_alg \; \label{interop-alg:border-msg-case:update-known-algs-2}

					\If{ timers[msg.sender\_alg] $\neq$ $\emptyset$ }{
						cancelTimer(timers[msg.sender\_alg]) \; \label{interop-alg:border-msg-case:rm-timers}
					}
					
				}
					
			}
		}
	}
}
\end{algorithm}

%%%% ALGO %%%%
The main idea is increase the set of forwarding nodes from only one running protocol with a set we called border nodes.
%
Intuition says that the bigger the size of this set the higher the reachability of nodes among broadcast sessions.
%
Border nodes are within the transmission range of other nodes running a different protocol (this is the case of $t_{4}, t_{5}, t_{6}$ in figure~\ref{packets-lost-ex}).


%%
Our interoperability mechanism works as follows.
%
Messages in the network are triplets $msg = (\tau, \rho, \varrho)$ where $\tau$ is either a control or broadcast message, $\rho$ refers to the running protocol at the sender and $\varrho$ is the message payload.
%
Control messages are beacons, which nodes exchanged periodically, required by some protocols to operate and broadcast messages contain the information to disseminate.
%
When nodes receive a message $msg$ that corresponds to the running protocol of the sender or the current receiver is a border node (line 10), $msg$ is broadcast; notice that the function \textit{broadcast()} (line~\ref{interop-alg:bro}) hides the concrete implementation of a \acl{bp}.
%
In case the running protocols between sender and receiver differs (line 13) we distinguish between two cases:

\begin{itemize}

\item \textbf{Control message reception.}
%
This case fosters the election of border nodes.
%
We keep track of senders from a foreign algorithm with the set \textit{known\_nodes} (line~\ref{interop-alg:save-sender}), potential candidates of border nodes are within this set.
%
If receivers previously hear a beacon from a sender running a known foreign protocol, $msg$ is ignored.
%
Otherwise (line: 16), we keep a reference of the knew known protocol (line~\ref{interop-alg:save-alg}) and an event is scheduled with a timer (line~\ref{interop-alg:new-timer}).
%
Once the timer expires (line~\ref{interop-alg:timer-expires}), we create a new border message \emph{b} (line~\ref{interop-alg:new-border-msg}) with the following fields:

\begin{itemize}

\item chosen border node (line~\ref{interop-alg:new-border-msg:chosen}). For the moment this is a random choice, we can think in other methods, like piggybacking senders remaining battery to control messages for selecting a border node with the higher value of energy. 
\item number of times $b$ is forward (line~\ref{interop-alg:new-border-msg:hops})
\item sender identifier (line~\ref{interop-alg:new-border-msg:sender})
\item sender running protocol (line~\ref{interop-alg:new-border-msg:sender-alg})
\item foreign protocol that fosters the border election procedure (line~\ref{interop-alg:new-border-msg:sender-alg})

\end{itemize}
%
Now we explore the case when nodes receive border messages (line 28).

\item \textbf{Border message reception.} Receivers of this message decrease by one the number of times \textit{msg} is forwarded (line:~\ref{interop-alg:border-msg-case:dec}).
%
In case this counter has remaining hops to perform, \textit{msg} is forward (lines: 30,~\ref{interop-alg:border-msg-case:fwd}).
%
Notice that the function \textit{forward()} simply performs a retransmission within the communication range of the receiver.
%
Given that all packets are identified in a unique way further copies of previously known packets are ignored.
%
Furthermore, only one out of two possible situations take place.

%%
Either the message receiver is a potential candidate of border node (line: 32) or the receiver belongs to those nodes who foster the selection procedure (line 36).
%
When the former situation occurs, the receiver is labeled as border node only if its identifier was chosen (line~\ref{interop-alg:border-msg-case:select-as-border}).

%
The latter situation serves to cancel ongoing timers that triggered the border selection procedure (line:~\ref{interop-alg:border-msg-case:rm-timers}) by other senders, that is, nodes running a foreign algorithm. 
%

\end{itemize}


\subsection{Evaluation scenario}

We simulate a situation where a crowd of humans are carrying mobile devices with WiFi capabilities (mobile phones, tablets, etc).
%
As shown in Figure~\ref{packets-lost-ex}, the area of communication is partitioned in two regions, dense and sparse.
%
The aim of having such an arrangement is to measure the impact of discovered border nodes with our interoperability mechanism that surrounds the dense region.
%
Both regions contain 400 nodes each and approximately the same surface where the dimensions of the whole area are 120 meters of height and width.
%
Nodes are spread uniformly in both regions and their initial position is chosen in a random way.
%
The number of nodes neighbors within the sparse region varies between 1 to 10 while those nodes in the dense region may have up to 40 neighbors.
%

\smallskip
\noindent
\textbf{Broadcast protocols.}
Running protocols on each node are defined in advance and nodes keep using this protocol until the end of the simulation.
%
Nodes in the sparse area use a controlled flooding approach that do not rely on an overlay.
%
When a node receives a broadcast message for the first time it starts a timer with a waiting period chosen randomly between the interval $[0,\dots,max]$. In our experiments the threshold \emph{max} is set to 10 ms.
%
If the number of received copies per message reaches the maximum number of allowed receptions before the timer expires, the message retransmission is canceled. Otherwise, the retransmission take place.
%
This approach was first propose in~\cite{Ni:1999:BSP:313451.313525} to deal with broadcast storms.
%%
Nodes in the dense area use an approximation of a connected-dominating set called multi-point relay technique (MPR)~\cite{qayyum:inria-00072756}.

\smallskip
\noindent
\textbf{Broadcast sessions.}
A total of 200 broadcast sessions are performed with a frequency of one second.
%
The source of each sessions is chosen randomly among all nodes.
%
The content of each message is randomly generated, unique and with a length of 32~Bytes. 


\smallskip
\noindent
\textbf{Simulator.}
We deploy our simulation with OMNET++4.6~\cite{Varga:2008:OOS:1416222.1416290} and the plug-in INET3.3 to model mobile ad hoc networks.
%
This framework provides a full implementation of the IEEE 802.11g standard, we use the standard model of energy consumption.
%

\smallskip
\noindent
\textbf{Transceiver settings.}
Nodes have a transmission range of 10 m.
%
This metric depends on the wireless communication chip used.
We model power consumption using the specifications of a commercial chip found in modern mobile devices.
%
We use Broadcom BCM4329 chip~\cite{BCM4329}, which supports Wi-Fi 802.11g on the 2.4 GHz frequency band.
%
Its energy consumption per operation modes is as follows: sleep (1mW), receive (10mW) and transmit (100mW).


\smallskip
\noindent
\textbf{Mobility.}
%\subsection*{Mobility}
We use the \emph{truncated-levy walk} (LW) technique to model mobility patterns of humans carrying wireless devices~\cite{Rhee:2011}.
%
The evaluation of our interoperability mechanism is performed in a first scenario with mobility and in a second scenario where nodes remain static.
%
In the scenario with mobility, nodes at the dense region move at a speed of 0.5 m/s and the speed of those at the sparse region is 2.0 m/s; these values are based on an estimation of the preferred walking speed of humans in~\cite{Browning390}.
%
The LW model generates traces of frequent short walks and occasional rides to distant locations.
%
Nodes move within their region of density in the communication area without getting out of their designated region.

	
\subsection{Preliminary Results}
\rcg{TBD...}

\begin{figure}[!t]
	\begin{center}
		\subfigure[Overall energy consumption]{
			\includegraphics[width=\linewidth]{energy_consump_static}
			\label{plots:energy-consump}
		}
		\subfigure[Relative error of received messages]{
			\includegraphics[width=\linewidth]{relative_error_rcv_msgs_static}
			\label{plots:re-rcv-msgs}
		}
				\subfigure[Network coverage]{
			\includegraphics[width=\linewidth]{netowrk_coverage_static}
			\label{plots:net-coverage}
		}
		\caption{Preliminary results of broadcast metrics in the static scenario.} 
		\label{fig:outbw}
	\end{center}
\end{figure}

\begin{figure}[!t]
  \centering
  \includegraphics[trim = 0 0 0 9mm, clip, scale=0.4]{fwd_types_dense_area_static}
  \caption{Forwarding sets grouped by type of retransmission}
  \label{plots:fwd-sets-type}
\end{figure}